% Please do not change %
\documentclass[11pt]{article} % font size
\usepackage[margin=1in]{geometry} % margins
\usepackage{amsmath,amsthm,amssymb} % for the  common "math symbols"
\usepackage{algorithm}
\usepackage{algpseudocode}
\usepackage[shortlabels]{enumitem} % for enumeration
\usepackage{booktabs} % if you need tables
\usepackage{listings}
\usepackage{color}
\usepackage{quiver}
\DeclareMathOperator*{\argmax}{argmax}
\DeclareMathOperator*{\argmin}{argmin}

% \theoremstyle{definition}
\newtheorem{theorem}{Theorem}[section]
\newtheorem{corollary}[theorem]{Corollary}
\newtheorem{lemma}[theorem]{Lemma}
\newtheorem{proposition}[theorem]{Proposition}
\newtheorem{conjecture}[theorem]{Conjecture}
\theoremstyle{definition}
\newtheorem{definition}[theorem]{Definition}
\theoremstyle{definition}
\newtheorem{fact}[theorem]{Fact}
\theoremstyle{definition}
\newtheorem{example}[theorem]{Example}

\lstset{
  frame=none,
  xleftmargin=2pt,
  stepnumber=1,
  numbers=left,
  numbersep=5pt,
  numberstyle=\ttfamily\tiny\color[gray]{0.3},
  belowcaptionskip=\bigskipamount,
  captionpos=b,
  escapeinside={*'}{'*},
  language=haskell,
  tabsize=2,
  emphstyle={\bf},
  commentstyle=\it,
  stringstyle=\mdseries\rmfamily,
  showspaces=false,
  keywordstyle=\bfseries\rmfamily,
  columns=flexible,
  basicstyle=\small\sffamily,
  showstringspaces=false,
  morecomment=[l]\%,
}
% Feel free to include additional packages (if needed), but make sure it does not override the margin/font requirements.
%
\setlength{\parindent}{0pt} % no indent for new paragraphs.
\setlength{\parskip}{5pt} % distance between paragraphs, which will be used when you have a blank line in your LaTeX code.


%%%%%%%%%%%%%%%%%%%%%%%%%%%%%%%%%%%%%%%%%%%%%%%%%%%%%%%%%%%%%%%%%%%%%%%%%%%%%%
\title{Categories, Proofs, and Processes: \\Assignment 5}
%
\author{Wira Azmoon Ahmad}
\date{26 Feb 2024}
%
\begin{document}
%
\maketitle

%You may use the usual \LaTeX\ environments to structure your answers:
%\begin{align}
%        \label{basegnn}
%        \mathbf{h}_u^{(t+1)} = \sigma \bigg( \mathbf{W}_\text{self}^{(t)} \mathbf{h}_u^{(t)}  + \mathbf{W}_\text{neigh}^{(t)} \sum_{v \in N(u)} \mathbf{h}_v^{(t)} + \mathbf{b}^{(t)} \bigg): 
%\end{align}


%\begin{table}[h]
%\caption{A simple table.}
%\centering
%\begin{tabular}[t]{lrrrr}
%\toprule
%Graphs & Can & Be  & Fun & Too \\ 
%\midrule 
%$G_1$ & $1$  & $2$ & $3$ & $4$ \\ 
%$G_2$ & $-1$  & $-2$ & $-3$ & $-4$ \\ 
%$G_3$ & $0.1$  & $0.2$ & $0.3$ & $0.4$ \\ 
%\bottomrule
%\end{tabular}
%\label{tab}
%\end{table}

% A functor $F:\mathcal{C} \rightarrow \mathcal{D}$ is
% \begin{itemize}
%     \item \textbf{faithful} if each map $F_{A,B}:\mathcal{C}(A,B) \rightarrow \mathcal{D}(F\ A, F\ B)$ is injective;
%     \item \textbf{full} if each map $F_{A,B}:\mathcal{C}(A,B) \rightarrow \mathcal{D}(F\ A, F\ B)$ is surjective
% \end{itemize}

\section*{Question 1}
\addtocounter{section}{1}
\subsection*{(a)}
The set-indexed product of a set of objects $\{A_i\}_{i\in I}$ is the Cartesian product
\[\prod_{i\in I}A_i = \left\{g: I \rightarrow \bigcup_{i \in I} A_i \middle| \forall i\in I: g(i) \in A_i\right\}\]
with its arrows $\pi_j(g) = g(j)$ for all $j\in I$.

Let $C$ be an object with a set of arrows $f_i:C\rightarrow A_i$.
Then the unique arrow $\langle f_i \rangle_{i\in I}$ is defined where for all $c\in C$,
\[\langle f_i \rangle_{i\in I}(c) := i \mapsto f_i(c)\]
For each $j\in I$, and for all $c\in C$,
\begin{equation}
\begin{aligned}
    \pi_j \circ \langle f_i \rangle_{i\in I} (c) &= \pi_j(i \mapsto f_i(c))\\
    &= (i \mapsto f_i(c))(j)\\
    &= f_j(c)
\end{aligned}
\end{equation}
so $\pi_j \circ \langle f_i \rangle_{i\in I} = f_j$.

Suppose there is another arrow $h: C\rightarrow\prod_{i\in I}A_i$ 
such that $\forall j\in I, \pi_j \circ h = f_j$.
Then $\forall c\in C$,
\begin{equation}
\begin{aligned}
    \pi_j \circ h (c) &= \pi_j(h(c))\\
    &= h(c)(j)\\
    &= f_j(c)\\
    &\implies h(c) :: j \mapsto f_j(c)
\end{aligned}
\end{equation}
So $h=\langle f_i \rangle_{i\in I}$ and the arrow is unique.

\subsection*{(b)}
When $|I|=2$, we get binary products.
The product of 2 objects $A_i,A_j$ for $i,j\in I$ is the object $A_i \times A_j$
with arrows $A_i\xleftarrow[]{\pi_i} A_i\times A_j \xrightarrow[]{\pi_j} A_j$
such that for any object $C$ and set of arrows $A_i\xleftarrow[]{\pi_i} C \xrightarrow[]{\pi_j} A_j$,
there is a unique arrow $\langle f_i, f_j\rangle:C\rightarrow A_i \times A_j$
such that $\pi_i \circ \langle f_i, f_j\rangle=f_i$ and $\pi_j \circ \langle f_i, f_j\rangle=f_j$.

When $|I|=1$, we get terminal objects.
The product of 1 object $A_i$ for $i\in I$ is the terminal object $A_i$
with the identity $1_{A_i}$
such that for any object $C$, the arrow $f_i:C\rightarrow A_i$,
is the unique arrow to $A_i$.

\subsection*{(c)}
Let $(L, \gamma)$ be the limit of $D:\mathcal{I}\rightarrow \mathcal{C}$.
Then the arrows are defined as 
\[\{\gamma_J : L\ \rightarrow D\ J\}_{J\in \text{Ob}(\mathcal{I})}\]
where the only arrows in $\mathcal{J}$ are identities $1_J$ so $D\ 1_J = 1_{D\ J}$.
Since $(L, \gamma)$ is the terminal object in \textbf{Cone}$(D)$,
for every other cone to $D, (A, \delta)$,
there is a unique arrow 
$g : A\rightarrow L$ such that $\forall I\in \text{Ob}(\mathcal{I}), \gamma_I \circ g = \delta_I$.

By using the functor $S: \textbf{Set} \rightarrow \textbf{Cat}$ from Assignment 3 Exercise 1,
let $S\ I = \mathcal{I}$, so there's an isomorphism $\phi: I \cong \text{Ob}(\mathcal{I})$.
Then set $\forall j\in I$,
\begin{itemize}
    \item $D\ \phi(j) := A_j$
    \item $L := \prod_{j\in I} A_j$
    \item $\gamma_{\phi(j)} := \pi_j$
    \item $A := C$
    \item $\delta_{\phi(j)} := f_j$
    \item $g:=\langle f_j\rangle_{j\in I}$
\end{itemize}
so the limit of the diagram matches the defintion of the set-indexed products.

\subsection*{(d)}
Suppose $\mathcal{D}$ has all set-indexed products.

Let $I$ be the set of indices, and $\{F_i:\mathcal{C}\rightarrow \mathcal{D}\}_{i\in I}$ 
be a set of functors.
First, we construct the product functor $\prod_{i\in I} F_i:\mathcal{C}\rightarrow \mathcal{D}$ 
and prove its functoriality.
\[\prod_{i\in I} F_i (C) := \prod_{i\in I} (F_i\ C)\]
\[\prod_{i\in I} F_i (f: C\rightarrow C') := \prod_{i\in I} (f_i: F_i\ C \rightarrow F_i\ C' )
:= \langle f_i \circ \pi_i \rangle_{i\in I}\]
where $\pi_i:\prod_{i\in I} F_i\rightarrow F_i\ C$, which exists in $\mathcal{D}$.

Let $C\xrightarrow[]{f} C'\xrightarrow[]{g} C''$ be arrows.
\begin{equation}
\begin{aligned}
    \prod_{i\in I} F_i (g \circ f) &= \langle (g\circ f)_i \circ \pi_i \rangle_{i\in I}\\
    &= \prod_{i\in I}(g\circ f)_i \circ\langle \pi_i \rangle_{i\in I}\\
    &= \prod_{i\in I}g_i \circ \prod_{i\in I}f_i \circ\langle \pi_i \rangle_{i\in I}\\
    &= \langle g_i \circ \pi_i \rangle_{i\in I} \circ \prod_{i\in I}f_i \circ\langle \pi_i \rangle_{i\in I}\\
    &= \langle g_i \circ \pi_i \rangle_{i\in I} \circ \langle f_i \circ \pi_i \rangle_{i\in I} \\
    &= \prod_{i\in I} F_i (g) \circ \prod_{i\in I} F_i (f)
\end{aligned}
\end{equation}
\begin{equation}
\begin{aligned}
    \prod_{i\in I} F_i (1_C) &= \langle 1_{F_i\ C} \circ \pi_i \rangle_{i\in I}\\
    &= \langle \pi_i \rangle_{i\in I}\\
    &= \prod_{i\in I} 1_{F_i\ C} \\
    &= 1_{\prod_{i\in I} (F_i\ C)} \\
    &= 1_{\prod_{i\in I} F_i (C)} \\
\end{aligned}
\end{equation}
So the product functor satisfies functoriality.

Next, we show the naturality of the functor projections.
Let $f: C\rightarrow C'$ be an arrow in $\mathcal{C}$, 
and $f_j: F_j\ C \rightarrow F_j\ C'$ be an arrow in $\mathcal{D}$.
For an arbitrary natural transformation $t_j: \prod_i F_i \rightarrow F_j$, 
define $t_{j,C} = \pi_j : \prod_i F_i\ C \rightarrow F_j\ C$,
and $t_{j,C'} = \pi'_j : \prod_i F_i\ C' \rightarrow F_j\ C'$,
which both exist as projections in $\mathcal{D}$.
Then by the property of the product, we have
\[f_j \circ \pi_j = \pi'_j \circ \langle f_i\circ \pi_i\rangle_i\]
So naturality is satisfied.

Let again $f: C\rightarrow C'$ be an arrow in $\mathcal{C}$, 
% and $f_j: F_j\ C \rightarrow F_j\ C'$ be an arrow in $\mathcal{D}$.
Suppose we have some functor $G$ and a set of natural transformations $s_i: G\rightarrow F_i$.
Define the natural transformation $\langle s_i \rangle_i: G\rightarrow \prod_i F_i$ such that
\[ \langle s_i \rangle_{i,C} = \langle s_{i,C} \rangle_i: G\ C \rightarrow \prod_i F_i\ C\]
By naturality of $s_i$, we have 
\[f_i \circ s_{i,C} = s_{i,C'}\circ G\ f\]
From the property of the product, we have that 
\[\pi_j \circ \langle s_i\rangle_{i, C} = s_{j, C},\quad
\pi'_j \circ \langle s_i\rangle_{i, C'} = s_{j, C'}\]
Then these imply that
\[f_j \circ \pi_j \circ \langle s_i\rangle_{i, C} = \pi'_j \circ \langle s_i\rangle_{i, C'} \circ G\ f\]
\[\pi'_j \circ \langle f_i\circ \pi_i\rangle_i \circ \langle s_i\rangle_{i, C} 
= \pi'_j \circ \langle s_i\rangle_{i, C'} \circ G\ f\]
Since this is true for all $i\in I$, then by the property of the product, 
\[\langle f_i\circ \pi_i\rangle_i \circ \langle s_i\rangle_{i, C} =\langle s_i\rangle_{i, C'} \circ G\ f\]
So the naturality of $\langle s_i \rangle_i$ is satisfied.
Since $\forall C$ in $\mathcal{C}$, and $\forall i\in I$, we have some $\pi_i$ that 
\[\pi_j \circ \langle s_i\rangle_{i, C} = s_{j, C}\]
where $t_j=\pi_j$, then
\[\forall j \in I: t_j \circ \langle s_i\rangle_i = s_j\]

Now, we show the uniqueness of the natural transformation. 
Suppose there is some other natural transformation $u:G\rightarrow \prod_i F_i$ 
such that $\forall j\in I,t_j \circ u = s_j$.
Then for all $C$ in $\mathcal{C}$, $\forall j\in I, \pi_j \circ u_C = s_{j,C}$ for some $\pi_j$.
Then there is a set of mappings 
\[\pi_j \circ u_c = s_{j,c} : G\ C \rightarrow F_j\ C\]
By property of the product, there is a unique mapping
\[\langle\pi_i \circ u_C\rangle_i = \langle s_{i,C}\rangle_i: G\ C\rightarrow \prod_i F_i\ C\]
such that $\forall j\in I, \pi_j \circ \langle\pi_i \circ u_C\rangle_i = s_{j,C}$.
Then for all $C$ in $\mathcal{C}$,
\begin{equation}
\begin{aligned}
    \langle\pi_i \circ u_C\rangle_i &= \langle s_{i,C}\rangle_i\\
    \langle\pi_i\rangle_i \circ u_C &= \langle s_i\rangle_{i,C}\\
    1_{\prod_i F_i\ C}\circ u_C &= \langle s_i\rangle_{i,C}\\
    u_C &= \langle s_i\rangle_{i,C}\\
    u &= \langle s_i\rangle_i
\end{aligned}
\end{equation}
Thus, $\langle s_i\rangle_i$ is the unique arrow.

\subsection*{(e)}

\section*{Question 2}
\addtocounter{section}{1}

\subsection*{(a)}
\subsection*{(b)}
\subsection*{(c)}
\subsection*{(d)}
\subsection*{(e)}
\subsection*{(f)}


\end{document}
